%%%%%%%%%%%%%%%%%%%%========== 見出しと番号 ==========%%%%%%%%%%%%%%%%%%%%

% 節の見出しに「§」，小節の見出しに「§§」を付ける．
% 後付け(\backmatter)では，main.texで空に再定義する．
\makeatletter
\renewcommand{\section}{%
	\if@slide\clearpage\fi
	\@startsection{section}{1}{\z@}%
	{\Cvs \@plus.5\Cdp \@minus.2\Cdp}% 前アキ
	{.5\Cvs \@plus.3\Cdp}% 後アキ
	{\normalfont\Large\headfont\presectionname\raggedright}%
}
\renewcommand{\subsection}{%
	\@startsection{subsection}{2}{\z@}%
	{\Cvs \@plus.5\Cdp \@minus.2\Cdp}% 前アキ
	{.5\Cvs \@plus.3\Cdp}% 後アキ
	{\normalfont\large\headfont\presubsectionname}%
}
\makeatother
\renewcommand{\presectionname}{\S}
\NewDocumentCommand{\presubsectionname}{}{\S\S}
\setcounter{tocdepth}{2}

% 式番号と脚注番号を，節ごとに振り直す．
\numberwithin{equation}{section}
\counterwithin*{footnote}{section}

%%%%%%%%%%%%%%%%%%%%========== キャプション ==========%%%%%%%%%%%%%%%%%%%%
\renewcommand{\thesubfigure}{\alph{subfigure}}

\makeatletter
\captionsetup{compatibility=false}
\captionsetup{%
	format = plain,
	labelsep = quad,
	font = small,
	width = .85\linewidth,
	subrefformat = parens,
	skip = 5\jsc@mpt
}
\captionsetup[subfigure]{%
	labelformat = parens,
	labelsep = space
}
\makeatother
